
%(BEGIN_QUESTION)
% Copyright 2014, Tony R. Kuphaldt, released under the Creative Commons Attribution License (v 1.0)
% This means you may do almost anything with this work of mine, so long as you give me proper credit

En interessant metode for elektrisk {\it varmekabel/tracing} er å føre svært høy strøm gjennom selve rørveggen, slik at metallrøret fungerer som en stor motstand. Varmen som utvikles i ``motstanden'' holder røret varmt nok til å hindre størkning av væsken inni.

\vskip 10pt

For å få så høy strøm brukes en nedtransformator som omformer industrispenning (typisk 240 VAC eller 480 VAC) til lavspenning med høy strøm for resistiv oppvarming av røret. Med 480 VAC enfase, bestem nødvendig nedtransformeringsforhold for å levere 1500 W varmeeffekt til et rør med ende-til-ende-resistans 0.75 ohm.

\vfil 

\underbar{file i00856}
\eject
%(END_QUESTION)





%(BEGIN_ANSWER)

Dette er en vurdert oppgave -- ingen svar eller hint oppgis!

%(END_ANSWER)





%(BEGIN_NOTES)

Effekt i en resistans er proporsjonal med kvadratet av spenningsfallet over resistansen:

$$P = {V^2 \over R}$$

Siden vi kjenner ønsket effekt (1500 W) og rørresistans (0.75 $\Omega$), kan vi løse for spenningsfallet over røret:

$$V = \sqrt{P R} = \sqrt{(1500 \hbox { W})(0.75 \> \Omega)} = 33.54 \hbox{ V}$$
 
Transformatoren må trappe ned 480 V til 33.54 V, og trenger derfor viklingsforhold:

$$\hbox{Turns ratio} = {480 \over 33.54} = 14.31$$

%INDEX% Electronics review: AC transformer circuit

%(END_NOTES)
